\VerdierPhysicalPageAnchor{0080}%
\VerdierSemanticAnchor{ch1-2-5-10}
\noindent\textbf{2.5.10.}\quad
\VerdierSemanticAnchor{ch1-2-5-10-multiadditive-setup}
Let
$F\colon\prod_{i\in[n]}C_i\longrightarrow C$ be a multiadditive functor.
Let us denote by

\VerdierSemanticAnchor{ch1-2-5-10-quotient-functors}
\begin{equation*}
  Q_i\colon\operatorname{comp}(C_i)\longrightarrow K(C_i)
\end{equation*}
the canonical functors. It follows from (2.5.6) and (2.5.7) that, when
countable products are representable in $C$, there exists one and only one
functor

\VerdierSemanticAnchor{ch1-2-5-10-pi-extension}
\begin{equation*}
  K_{\Pi}(F)\colon
  \prod_{i\in[n]}K(C_i)\longrightarrow K(C)
\end{equation*}
such that the following diagram is commutative:

\VerdierSemanticAnchor{ch1-2-5-10-eq1}
\begin{equation}
  \resizebox{0.84\textwidth}{!}{%
    $\begin{tikzcd}[ampersand replacement=\&,column sep=6.8em,row sep=5.0em]
      \displaystyle\prod_{i\in[n]}\operatorname{comp}(C_i)
        \arrow[r,"{\operatorname{comp}_{\Pi}F}"]
        \arrow[d,"{\prod_{i\in[n]}Q_i}"']
        \& \operatorname{comp}(C)
          \arrow[d,"Q"] \\
      \displaystyle\prod_{i\in[n]}K(C_i)
        \arrow[r,"{K_{\Pi}F}"]
        \& K(C)\mathpunct{.}
    \end{tikzcd}$%
  }
  \tag{2.5.10.1}
\end{equation}

\VerdierSemanticAnchor{ch1-2-5-10-twisted-structure}
Moreover, the functor $\prod_{i\in[n]}Q_i$ being a strict
$\mathbb Z[n]$-functor, there exists on the functor $K_{\Pi}F$ one and only
one structure of graded $\mathbb Z[n]$-functor twisted by the Koszul cocycle
relative to the ordered basis $e_1,\ldots,e_n$ of $\mathbb Z[n]$, such that
one has the equality between twisted $\mathbb Z[n]$-functors:

\VerdierSemanticAnchor{ch1-2-5-10-eq2}
\begin{equation}
  K_{\Pi}F\prod_{i\in[n]}Q_i
  =Q\operatorname{comp}_{\Pi}F.
  \tag{2.5.10.2}
\end{equation}

\VerdierSemanticAnchor{ch1-2-5-10-morphism-setup}
Finally, let $F$ and $F'$ be two multiadditive functors and let
$m\colon F\longrightarrow F'$ be a morphism of functors. It follows from
(2.5.7) that there exists one and only one morphism of functors:

\VerdierSemanticAnchor{ch1-2-5-10-eq3}
\begin{equation}
  K_{\Pi}m\colon K_{\Pi}F\longrightarrow K_{\Pi}F'
  \tag{2.5.10.3}
\end{equation}
such that:

\VerdierSemanticAnchor{ch1-2-5-10-eq4}
\begin{equation}
  K_{\Pi}m\star\prod_{i\in[n]}Q_i
  =Q\star\operatorname{comp}_{\Pi}m.
  \tag{2.5.10.4}
\end{equation}

\VerdierSemanticAnchor{ch1-2-5-10-pi-terminology}
The functor $K_{\Pi}F$ is called the $\Pi$-\emph{extension} of the functor
$F$ to complexes up to homotopy. The morphism $K_{\Pi}m$ is called the
$\Pi$-\emph{extension} of the morphism $m$ to complexes up to homotopy.

\VerdierSemanticAnchor{ch1-2-5-10-sigma-analogue}
When countable direct sums are representable in $C$, one defines analogously
the functor $K_{\Sigma}F$: the $\Sigma$-\emph{extension} of the functor $F$
to complexes up to homotopy, and the morphism $K_{\Sigma}m$: the
$\Sigma$-\emph{extension} of the morphism $m$ to complexes up to homotopy.
